\section{Impacto en el empleo por niveles: qué puestos cambian primero}

La deflación del trabajo de cuello blanco no ocurre de manera uniforme. Coding/Agent afectará primero a los puestos cuyas tareas tienen un alto grado de digitalización, cuyos resultados son verificables y cuyos procesos pueden descomponerse. No necesariamente eliminará de inmediato una profesión completa, pero primero reorganizará la proporción de tareas dentro de cada puesto.

\subsection{Matriz de impacto por puesto}

\noindent\begin{longtable}{p{0.22\textwidth}p{0.34\textwidth}p{0.35\textwidth}}
\toprule
Puesto/tarea & Parte que se comprime primero & Parte más difícil de sustituir \\
\midrule
Programador júnior & Código repetitivo, bugs simples, scripts, encapsulamiento de interfaces & Criterio de arquitectura, comprensión de sistemas complejos, responsabilidad por la puesta en producción. \\
Análisis de datos & Limpieza, SQL, reportes, primeras versiones de visualizaciones & Definición de métricas, interpretación del negocio, juicio causal. \\
Producto/operaciones & Documentación, resúmenes de la competencia, configuración de procesos & Comprensión profunda de los usuarios, toma de decisiones entre alternativas, coordinación organizacional. \\
Asistente de investigación & Organización de bibliografía, scripts de experimentos, síntesis de resultados & Planteamiento de preguntas, diseño experimental, explicación de los fracasos. \\
Gestión y coordinación & Sincronización del estado, minutas de reuniones, calendarización & Resolución de conflictos, decisiones sobre recursos, asunción de responsabilidades. \\
\bottomrule
\end{longtable}

\begin{importantbox}{La forma real del cambio en las profesiones}
La IA no necesariamente sustituye un puesto de una sola vez, sino que primero sustituye tareas. Los puestos se reorganizarán: las partes de bajo criterio, alta repetición y resultados verificables se automatizan; las partes de alta responsabilidad, mucho contexto y decisiones difíciles entre alternativas cobran más importancia.
\end{importantbox}

\subsection{Resumen de la sección}

El impacto en el empleo debe analizarse con la granularidad de las tareas. Lo primero que se comprime son las tareas digitalizadas, repetitivas y verificables; lo más difícil de sustituir son la responsabilidad, el criterio, la organización y las decisiones sobre valores.
